\subsection{半单李代数相对于半单作用的不变量}

在本节中，假定 k 的特征为零。

命题 13. 设 g 是一个半单李代数，a 是 g 的一个在 g 中约化的子代数，m 是 a 在 g 中的交换子。子代数 m 满足第 3 节引理 2 的条件 (1) 与 (2)；它在 g 中是约化的。

把第一章，§3，第 5 节的命题 6 应用于 a-模 g，得 \(g = m \oplus [a, g]\) 。设 B 是 g 的 Killing 型，并设 \(x \in a, y \in m, z \in g\) 。那么，

\[
\mathrm{B} ([ z, x ], y) = \mathrm{B} (z, [ x, y ]) = 0 \quad \text { since } [ x, y ] = 0,
\]

这表明 m 关于 B 与 \([a,g]\) 正交。由于 B 非退化，且 \(g = m \oplus [a,g]\) ，这蕴含 B 在 m 上的限制非退化；于是引理 2 的条件 (1) 满足。

现在设 \(x \in m\) ，设 s 与 n 是它的半单分量与幂零分量。ad x 的半单分量是 ad s，参看~第一章，§6，第 3 节。由于 ad x 在 a 上为零，故由命题 4 (i) 知 ad s 也在 a 上为零。于是 \(s \in m\) ，从而 \(n = x - s \in m\) ，引理 2 的条件 (2) 满足。

注. \(\mathfrak{m}\) 在 \(\mathfrak{g}\) 中的交换子不一定化为 \(\mathfrak{a}\) ，参看~习题 4。

命题 14. 设 g 是一个半单李代数，A 是一个群，r 是从 A 到 \(\text{Aut}(\mathfrak{g})\) 的一个同态。设 m 是由在 \(r(A)\) 下不变的元素组成的 g 的子代数。假定线性表示 r 是半单的。那么 m 满足第 3 节引理 2 的条件 (1) 与 (2)；它在 g 中是约化的。

证明与前一命题的证明类似：

设 \(\mathfrak{g}^{+}\) 是 \(\mathfrak{g}\) 的由诸 \(r(a)x - x\) （\(a \in \mathrm{A}\) ，\(x \in \mathfrak{g}\) ）生成的向量子空间。向量空间 \(\mathfrak{g}' = \mathfrak{m} + \mathfrak{g}^{+}\) 在 \(r(\mathrm{A})\) 下稳定。设 \(\mathfrak{n}\) 是 \(\mathfrak{g}'\) 在 \(\mathfrak{g}\) 中的余子空间，且在 \(r(\mathrm{A})\) 下稳定。若 \(x \in \mathfrak{n}\) ，\(a \in \mathrm{A}\) ，则 \(r(a)x - x \in \mathfrak{n} \cap \mathfrak{g}^{+} = 0\) ，故 \(x \in \mathfrak{m}\) ，从而 \(x = 0\) ，由于 \(\mathfrak{m} \cap \mathfrak{n} = 0\) 。于是 \(\mathfrak{g} = \mathfrak{g}' = \mathfrak{m} + \mathfrak{g}^{+}\) 。设 B 是 \(\mathfrak{g}\) 的 Killing 型，并设 \(y \in \mathfrak{m}\) ，\(a \in \mathrm{A}\) ，\(x \in \mathfrak{g}\) 。那么

\[
\begin{array}{r l} \mathrm{B} (y, r (a) x - x) & = \mathrm{B} (y, r (a) x) - \mathrm{B} (y, x) \\ & = \mathrm{B} (r (a ^ {- 1}) y, x) - \mathrm{B} (y, x) \\ & = \mathrm{B} (y, x) - \mathrm{B} (y, x) = 0. \end{array}
\]

于是 m 与 \(g^{+}\) 关于 B 正交。由此得出 B 在 m 上的限制非退化；从而得到引理 2 的条件 (1)。条件 (2) 由结构转移立即可得。

